% Ablauf Train -> Play -> Sim2Real für Lauf-Policies am Beispiel unitree_rl_mjlab.
% Eigene Darstellung in Anlehnung an UnitreeRLMjlab (README, S. 1: "Train → Play → Sim2Real";
% Belege UnitreeRLMjlab.02, .04, .06, .07).
% Stil: ba-grafik.tex. Einbinden: % Ablauf Train -> Play -> Sim2Real für Lauf-Policies am Beispiel unitree_rl_mjlab.
% Eigene Darstellung in Anlehnung an UnitreeRLMjlab (README, S. 1: "Train → Play → Sim2Real";
% Belege UnitreeRLMjlab.02, .04, .06, .07).
% Stil: ba-grafik.tex. Einbinden: % Ablauf Train -> Play -> Sim2Real für Lauf-Policies am Beispiel unitree_rl_mjlab.
% Eigene Darstellung in Anlehnung an UnitreeRLMjlab (README, S. 1: "Train → Play → Sim2Real";
% Belege UnitreeRLMjlab.02, .04, .06, .07).
% Stil: ba-grafik.tex. Einbinden: % Ablauf Train -> Play -> Sim2Real für Lauf-Policies am Beispiel unitree_rl_mjlab.
% Eigene Darstellung in Anlehnung an UnitreeRLMjlab (README, S. 1: "Train → Play → Sim2Real";
% Belege UnitreeRLMjlab.02, .04, .06, .07).
% Stil: ba-grafik.tex. Einbinden: \input{Bilder/tikz/rl-ablauf} in einer figure-Umgebung.
\begin{tikzpicture}[ba, node distance=8mm]
  \tikzset{blk/.style={ba block, text width=40mm, minimum height=16mm}}
  \node[blk, ba block hervor] (T) {\textbf{Train}\\Agent interagiert mit der MuJoCo-Simulation};
  \node[blk, right=10mm of T] (P) {\textbf{Play}\\Wiedergabe zur Prüfung des Verhaltens};
  \node[blk, right=10mm of P] (S) {\textbf{Sim2Real}\\Einsatz am\\realen Roboter};
  \draw[ba pfeil] (T) -- (P);
  \draw[ba pfeil] (P) -- (S);
  % Lernschleife unter Train
  \node[ba block, minimum width=15mm] (A) at ($(T.south)+(-1.6,-1.9)$) {Agent};
  \node[ba block, minimum width=15mm] (M) at ($(T.south)+(1.6,-1.9)$) {MuJoCo};
  \draw[ba pfeil] (A.north) to[bend left=30] node[above, font=\footnotesize] {Aktion} (M.north);
  \draw[ba pfeil] (M.south) to[bend left=30] node[below, font=\footnotesize] {Belohnung} (A.south);
  \node[ba gruppe, fit=(A)(M), inner sep=6pt, yshift=0pt] (G) {};
  \draw[ba hilfslinie] (T.south) -- (G.north);
  % ONNX-Export
  \draw[ba pfeil, dash pattern=on 3pt off 2pt] (G.east) -| node[pos=0.25, below, font=\footnotesize] {Export als ONNX} (S.south);
\end{tikzpicture}
 in einer figure-Umgebung.
\begin{tikzpicture}[ba, node distance=8mm]
  \tikzset{blk/.style={ba block, text width=40mm, minimum height=16mm}}
  \node[blk, ba block hervor] (T) {\textbf{Train}\\Agent interagiert mit der MuJoCo-Simulation};
  \node[blk, right=10mm of T] (P) {\textbf{Play}\\Wiedergabe zur Prüfung des Verhaltens};
  \node[blk, right=10mm of P] (S) {\textbf{Sim2Real}\\Einsatz am\\realen Roboter};
  \draw[ba pfeil] (T) -- (P);
  \draw[ba pfeil] (P) -- (S);
  % Lernschleife unter Train
  \node[ba block, minimum width=15mm] (A) at ($(T.south)+(-1.6,-1.9)$) {Agent};
  \node[ba block, minimum width=15mm] (M) at ($(T.south)+(1.6,-1.9)$) {MuJoCo};
  \draw[ba pfeil] (A.north) to[bend left=30] node[above, font=\footnotesize] {Aktion} (M.north);
  \draw[ba pfeil] (M.south) to[bend left=30] node[below, font=\footnotesize] {Belohnung} (A.south);
  \node[ba gruppe, fit=(A)(M), inner sep=6pt, yshift=0pt] (G) {};
  \draw[ba hilfslinie] (T.south) -- (G.north);
  % ONNX-Export
  \draw[ba pfeil, dash pattern=on 3pt off 2pt] (G.east) -| node[pos=0.25, below, font=\footnotesize] {Export als ONNX} (S.south);
\end{tikzpicture}
 in einer figure-Umgebung.
\begin{tikzpicture}[ba, node distance=8mm]
  \tikzset{blk/.style={ba block, text width=40mm, minimum height=16mm}}
  \node[blk, ba block hervor] (T) {\textbf{Train}\\Agent interagiert mit der MuJoCo-Simulation};
  \node[blk, right=10mm of T] (P) {\textbf{Play}\\Wiedergabe zur Prüfung des Verhaltens};
  \node[blk, right=10mm of P] (S) {\textbf{Sim2Real}\\Einsatz am\\realen Roboter};
  \draw[ba pfeil] (T) -- (P);
  \draw[ba pfeil] (P) -- (S);
  % Lernschleife unter Train
  \node[ba block, minimum width=15mm] (A) at ($(T.south)+(-1.6,-1.9)$) {Agent};
  \node[ba block, minimum width=15mm] (M) at ($(T.south)+(1.6,-1.9)$) {MuJoCo};
  \draw[ba pfeil] (A.north) to[bend left=30] node[above, font=\footnotesize] {Aktion} (M.north);
  \draw[ba pfeil] (M.south) to[bend left=30] node[below, font=\footnotesize] {Belohnung} (A.south);
  \node[ba gruppe, fit=(A)(M), inner sep=6pt, yshift=0pt] (G) {};
  \draw[ba hilfslinie] (T.south) -- (G.north);
  % ONNX-Export
  \draw[ba pfeil, dash pattern=on 3pt off 2pt] (G.east) -| node[pos=0.25, below, font=\footnotesize] {Export als ONNX} (S.south);
\end{tikzpicture}
 in einer figure-Umgebung.
\begin{tikzpicture}[ba, node distance=8mm]
  \tikzset{blk/.style={ba block, text width=40mm, minimum height=16mm}}
  \node[blk, ba block hervor] (T) {\textbf{Train}\\Agent interagiert mit der MuJoCo-Simulation};
  \node[blk, right=10mm of T] (P) {\textbf{Play}\\Wiedergabe zur Prüfung des Verhaltens};
  \node[blk, right=10mm of P] (S) {\textbf{Sim2Real}\\Einsatz am\\realen Roboter};
  \draw[ba pfeil] (T) -- (P);
  \draw[ba pfeil] (P) -- (S);
  % Lernschleife unter Train
  \node[ba block, minimum width=15mm] (A) at ($(T.south)+(-1.6,-1.9)$) {Agent};
  \node[ba block, minimum width=15mm] (M) at ($(T.south)+(1.6,-1.9)$) {MuJoCo};
  \draw[ba pfeil] (A.north) to[bend left=30] node[above, font=\footnotesize] {Aktion} (M.north);
  \draw[ba pfeil] (M.south) to[bend left=30] node[below, font=\footnotesize] {Belohnung} (A.south);
  \node[ba gruppe, fit=(A)(M), inner sep=6pt, yshift=0pt] (G) {};
  \draw[ba hilfslinie] (T.south) -- (G.north);
  % ONNX-Export
  \draw[ba pfeil, dash pattern=on 3pt off 2pt] (G.east) -| node[pos=0.25, below, font=\footnotesize] {Export als ONNX} (S.south);
\end{tikzpicture}
